\documentclass[a4paper,11pt]{article}
\usepackage[margin=1cm]{geometry}
\usepackage[utf8]{inputenc}
\usepackage[spanish]{babel}
\usepackage{amsmath}
\usepackage{enumitem}

\pagestyle{empty}
\setlist[enumerate]{nosep, leftmargin=1.5em}

\newcommand{\prob}[2]{\item #1 \textbf{(#2 pts)}}

\begin{document}
\noindent
\begin{minipage}[t][0.47\textheight][t]{0.48\textwidth}
  \begin{center}
    \textbf{\large Mecánica de los Fluídos}\\
    \textbf{\Large Evaluación - Continuidad y Bernoulli}
  \end{center}
  \begin{enumerate}
    \prob{Una cañería de $d_1=1$ m se reduce a $d_2=0,5$ m. Si $v_1=10$ m/s, calcule $v_2$.}{2}
    \prob{La cañería forzada de una central tiene $d_1=2$ m con $v_1=2$ m/s. Al llegar a la turbina $d_2=0,5$ m. Calcule la velocidad de impacto.}{2}
    \prob{Un tanque cilíndrico $r=50$ cm contiene $m=10000$ kg de agua. Si se perfora un orificio en la base, encuentre $v_2$.}{1,5}
    \prob{¿A qué velocidad sale un líquido por un orificio a 3 m de profundidad?}{1}
    \prob{Un río tiene $Q=11$ m$^3$/s, ancho $l=10$ m y profundidad $h=2$ m. Calcule la velocidad y el aumento si la sección se reduce un 30\%.}{1,5}
    \prob{Tubo de Pitot: el agua llega a $H=0,20$ m. ¿A qué velocidad va la corriente?}{1}
    \prob{Agua va de $s_1$ a $s_2$: $d_1=25$ mm, $p_1=345$ kPa, $v_1=3$ m/s. $s_2$: $d_2=50$ mm, 2 m arriba. Calcule $p_2$ (sin pérdidas).}{2}
  \end{enumerate}
\end{minipage}
\hfill
\begin{minipage}[t][0.47\textheight][t]{0.48\textwidth}
  \begin{center}
    \textbf{\large Mecánica de los Fluídos}\\
    \textbf{\Large Evaluación - Continuidad y Bernoulli}
  \end{center}
  \begin{enumerate}
    \prob{Un caudal de 20 L/s sale por un conducto de 0,1 m de diámetro y una boquilla de 5 cm. Determine las velocidades en ambas secciones.}{2}
    \prob{Un tanque cilíndrico de 2 m de diámetro se vacía en 10 s por un tubo de 0,5 m de diámetro. Calcule la velocidad de salida.}{1,5}
    \prob{¿Con qué velocidad sale un líquido por un orificio que se encuentra a una profundidad de 3 m?}{1}
    \prob{Un tanque de sección cuadrada $l=50$ cm contiene $m=31000$ kg. Si se perfora un orificio en la base, calcule $v_2$.}{2}
    \prob{Un biodigestor está lleno hasta $h=12$ m. Si se abre un orificio en la base, determine la velocidad de salida.}{1,5}
    \prob{Un recipiente cónico invertido contiene agua hasta $h=8$ m. Con orificio en la base, determine $v(t)$.}{2}
    \prob{Un conducto en V tiene sección superior de 2 m$^2$ y inferior de 0,5 m$^2$. Caudal de 1 m$^3$/s en la parte superior. Calcule las velocidades inicial y final.}{1}
  \end{enumerate}
\end{minipage}

\vspace{0.2cm}

\noindent
\begin{minipage}[t][0.47\textheight][t]{0.48\textwidth}
  \begin{center}
    \textbf{\large Mecánica de los Fluídos}\\
    \textbf{\Large Evaluación - Continuidad y Bernoulli}
  \end{center}
  \begin{enumerate}
    \prob{Un sistema de tuberías consta de tres tubos en serie: diámetros 10 cm, 8 cm y 6 cm. Caudal de 0.5 m$^3$/s. Calcule la velocidad en cada parte.}{2}
    \prob{Se desea que un estanque posea circulación de agua. Orificio $r=5$ mm bajo $H=0,5$ m. Determine el caudal de recirculación $Q=S \cdot v$.}{1,5}
    \prob{Un tanque cilíndrico de 1 m de diámetro se llena en 5 min por un tubo de 0,2 m de diámetro. Calcule la velocidad de entrada.}{1}
    \prob{Un canal rectangular: ancho 3 m, altura 2 m, caudal 4 m$^3$/s. Calcule la velocidad del agua.}{1}
    \prob{Un balde de 5 L se llena en 1 min con agua de una canilla. Si se reduce la mitad la sección, calcule $v_2$.}{1,5}
    \prob{Un recipiente cónico: diámetro base 4 m, altura 6 m. Caudal constante de 0,2 m$^3$/s. Calcule la velocidad en la parte más estrecha.}{2}
    \prob{Un depósito prisma rectangular: $w=6$ m, $l=8$ m, $h=10$ m. Orificio en base de 1 cm. Calcule la velocidad de descenso del agua al inicio.}{1}
  \end{enumerate}
\end{minipage}
\hfill
\begin{minipage}[t][0.47\textheight][t]{0.48\textwidth}
  \begin{center}
    \textbf{\large Mecánica de los Fluídos}\\
    \textbf{\Large Evaluación - Continuidad y Bernoulli}
  \end{center}
  \begin{enumerate}
    \prob{Un recipiente irregular contiene agua hasta $h=15$ m. Con orificio pequeño, determine la velocidad de salida.}{2}
    \prob{Un tanque tiene forma de esfera y contiene agua hasta $h=4$ m. Con orificio en la base, determine la ecuación de $v$ en función del radio y el tiempo.}{1,5}
    \prob{Determinar a qué altura se debe abrir un orificio para que el líquido salga con $v=9,8$ m/s.}{1}
    \prob{Un tanque de almacenamiento en forma de prisma rectangular: base 4 m x 6 m, altura 3 m. Se llena con agua a 0.1 m$^3$/s por una tubería de 0.2 m de diámetro. Calcule el tiempo de llenado.}{1,5}
    \prob{Un conducto con forma de V tiene sección transversal de 2 m$^2$ en la parte superior y 0.5 m$^2$ en la inferior. Si se bombea agua a 1 m$^3$/s en la parte superior, calcule las velocidades inicial y final.}{1}
    \prob{Un tanque cilíndrico de 1 m de diámetro se llena con agua a través de un tubo de 0.2 m de diámetro en 5 minutos. Calcula la velocidad de entrada.}{1}
    \prob{Un canal rectangular tiene un ancho de 3 metros y una altura de 2 metros. Si se envía agua a un caudal de 4 metros cúbicos por segundo, calcula la velocidad del agua en el canal.}{2}
  \end{enumerate}
\end{minipage}

\end{document}